\clearpage
\section{Allocation and beta opening of the original private-cell program}
\label{sec:proposition-opening}
This milestone proves a concrete contextual equivalence starting from the
unchanged left-hand Proposition~12.2 translation, first allocating its private
store and then opening its outer let-body by micro-beta. The theorem holds
in arbitrary finite variable and atom environments and in every original
binding-free closing context at unrestricted surrounding depth. It has no
unproved determinism, robustness or translation-correctness premise.

The endpoint is a specific graph containing the actual opened translated
body and the retained allocated store. It is not the right-hand program.
The checked reference resolves to this store in exactly two atom hops, and
the store is still attached to the original numeral one. The private state
has not yet been eliminated.

\paragraph{Relationship to the paper.}
The proof uses the reference-creation and micro-beta templates from the
\href{https://arxiv.org/html/1907.01257v5\#S12.SS3}{paper's Figure~22 chain}.
The paper first normalizes the internal sharing and box arrangement. Here
allocation and outer beta are applied directly to the unchanged translation,
after only representation isomorphisms removing genuinely empty forests and
reassociating graph assembly. Thus the opened intermediate still contains
the original internal distributors and is not asserted to be literally the
fourth canonical drawing of Figure~22. Structural normalization to the
parametricity instance, the parametricity link, and the final normalization
remain. This ordering changes neither a term nor a translation clause, and
every link used here has its own proved contextual consequence.

\paragraph{The concrete chain.}
Let $B$ be the actual translation of the outer let-body, let $N$ be the
numeral-one corolla, and let $i$ be the primitive instance edge. Write
$R=(C^\diamond_1;\mathsf{store});N$ for the retained store suffix, where $C^\diamond_1$ is the canonical
one-use atom tree. The proved chain is
\[
 \begin{aligned}
 P_{0,0}&\cong\mathsf{before}_{\beta}(B,\mathsf{before}_{\mathrm{ref}}(N))\\
 &\simeq\mathsf{before}_{\beta}(B,\mathsf{after}_{\mathrm{ref}}(N))\\
 &\cong\mathsf{before}_{\beta}(B,i);R\\
 &\simeq\mathsf{after}_{\beta}(B,i);R=O.
 \end{aligned}
\]
The isomorphisms fix the external interface. The contextual equivalences use
the universal cost relation; the local directed improvement results are
retained separately. Crucially, beta's stable argument is $i$, not the whole
instance--sharing-tree--store graph. Stability of the latter is explicitly
refuted in the checked development.

\Needspace{12\baselineskip}
\subsection{Only genuinely empty forests are removed}

A graph with no outputs can be followed by an edge-free graph of interface
$\epsilon\Rightarrow\epsilon$ without changing its recursive isomorphism
class. The old left occurrence fixes inputs, has no output equations to
check, and covers every target edge because a right-component edge is
impossible. The finite edge-cover theorem constructs the isomorphism.
Tensoring an arbitrary graph with an edge-free empty-interface component
on either side is handled similarly. On the right the original append-empty
interface equalities are transported explicitly. These lemmas do not discard
any nonempty weakening forest.

\lean{Copying/ClosedTranslationUnits}{10}{17}{ClosedTranslation.serial_empty}

\lean{Copying/ClosedTranslationUnits}{18}{27}{ClosedTranslation.tensor_empty_right}

\lean{Copying/ClosedTranslationUnits}{28}{37}{ClosedTranslation.tensor_empty_left}

A distributor with zero resources has no edges: any edge would contain an
element of $\mathrm{Fin}(0)$. This proves emptiness for either sharing sort,
then for the paired zero-environment distributor and its single-use interface
cast. Consequently a closed constant translation is isomorphic to its literal
constant corolla. A closed unary translation loses only its empty trailing
distributor, and a closed abstraction loses only its empty outer distributor;
its original box body is retained unchanged.

\lean{Copying/ClosedTranslationUnits}{38}{41}{ClosedTranslation.distributor_empty}

\lean{Copying/ClosedTranslationUnits}{42}{45}{ClosedTranslation.environment_empty}

\lean{Copying/ClosedTranslationUnits}{46}{48}{ClosedTranslation.singleUse_empty}

\lean{Copying/ClosedTranslationUnits}{49}{57}{ClosedTranslation.constant}

\lean{Copying/ClosedTranslationUnits}{58}{61}{ClosedTranslation.unary}

\lean{Copying/ClosedTranslationUnits}{62}{64}{ClosedTranslation.abstraction}

The closed binary operand pair has no output ports to shuffle. Join the two
original operand occurrences with the tensor boundary; input zero and input
one are checked separately, output preservation is vacuous, and the two edge
tags exhaust the target. The resulting isomorphism identifies the actual
Figure~16 pair with the tensor of its operands. Closed application then
removes its empty final distributor and uses this pair isomorphism under
the unchanged application head.

\lean{Copying/ClosedTranslationUnits}{65}{85}{ClosedTranslation.pair}

\lean{Copying/ClosedTranslationUnits}{86}{90}{ClosedTranslation.application}

\Needspace{12\baselineskip}
\subsection{Moving a closed-output branch past a neighbouring connection}

For $G:I\Rightarrow\epsilon$, $W:K\Rightarrow J$ and $R:J\Rightarrow O$,
the two graphs $(G\otimes W);R$ and $G\otimes(W;R)$ have the same ordered
interface $IK\Rightarrow O$. Unlike the earlier nullary-neighbour lemma,
$W$ may have arbitrary inputs; the instance edge needs this generality.
These two expressions define the left and right assembly shapes.

\lean{Copying/ClosedBranchSerial}{11}{11}{ClosedBranchSerial.left}

\lean{Copying/ClosedBranchSerial}{12}{98}{ClosedBranchSerial.right}

The isomorphism reassociates the three edge tags and the three component
internal-vertex tags, retaining the junction indexed by $J$. Its inverse
restores the old tags, and the two inverse equations follow by tag cases.
For ports of $G$, output cases are impossible; its inputs keep their original
positions. Inputs of $W$ retain the second input block, and its outputs become
the same serial junctions as inputs of $R$. Outputs of $R$ keep their exterior
positions. These three component port equations prove both ordered incidence
laws by mapping each source or target list. Labels are unchanged, including
all nested box bodies. This is a representation isomorphism, not a rewrite
that moves state across a box.

\lean{Copying/ClosedBranchSerial}{99}{134}{ClosedBranchSerial.iso}

\Needspace{12\baselineskip}
\subsection{Literal program geometry and the stable beta fragment}

\termdef{private-cell-opened-body}
The body term is exactly
$\lambda f.(\lambda w.!x)(f\ ())$ in the single-variable environment for $x$.
Under its sequencing binder $x$ is index two. Its translation $B$ has depth
two and interface $*\Rightarrow *$. The numeral graph is the original
constant-one corolla at the requested allowance. These definitions merely
name parts of the original desugared term.

\lean{Copying/PropositionAllocationGeometry}{12}{14}{Proposition12_2.bodyTerm}

\lean{Copying/PropositionAllocationGeometry}{15}{16}{Proposition12_2.body}

\lean{Copying/PropositionAllocationGeometry}{17}{18}{Proposition12_2.numeral}

Padding a primitive corolla to a larger allowance gives that same primitive
corolla at the new allowance. Induct on the inequality between allowances:
the reflexive case uses self-padding, and a successor uses the original
padding-successor equation and primitive-raising equation. Combining this
with the closed-constant isomorphism gives the padded constant translation.
The source isomorphism unfolds only the outer original application clause,
uses self-padding of the depth-three abstraction, commutes padding through
the reference unary clause, and applies the closed constructor isomorphisms.
The entire nested body remains $B$.

\lean{Copying/PropositionAllocationGeometry}{19}{24}{Proposition12_2.padTo_primitive}

\lean{Copying/PropositionAllocationGeometry}{25}{32}{Proposition12_2.constant_padding}

\lean{Copying/PropositionAllocationGeometry}{33}{44}{Proposition12_2.source_iso}

\termdef{private-cell-store-suffix}
The store suffix $R$ is the serial composite of the actual canonical
one-use atom tree, the store edge and the numeral-one graph. The allocated
graph retains the original application/closure with the freshly allocated
argument. The beta-source graph exposes only the instance as the argument
and places $R$ outside it. The opened graph replaces this beta-source fragment
by the actual micro-beta reduct, keeping the same $R$. All three are explicit
depth-three closed graph expressions.

\lean{Copying/PropositionAllocationGeometry}{45}{47}{Proposition12_2.storeSuffix}

\lean{Copying/PropositionAllocationGeometry}{48}{50}{Proposition12_2.allocated}

\lean{Copying/PropositionAllocationGeometry}{51}{54}{Proposition12_2.betaSource}

\lean{Copying/PropositionAllocationGeometry}{55}{57}{Proposition12_2.opened}

Three serial associativity isomorphisms factor the allocated argument into
instance followed by $R$. Tensor congruence retains the closed closure, the
closed-output branch isomorphism moves $R$ to the outside of that tensor,
and inverse serial association exposes the application/closure/instance
beta fragment. No store edge is deleted, duplicated or included in beta's
stable argument.

\lean{Copying/PropositionAllocationGeometry}{58}{66}{Proposition12_2.allocated_argument_iso}

\lean{Copying/PropositionAllocationGeometry}{67}{72}{Proposition12_2.allocated_iso}

\Needspace{12\baselineskip}
\subsection{The actual binding-free observation contexts}

The allocation context applies a fixed closed closure to a single closed
term hole. Its address follows the right tensor component and then the
right side of the application-head serial composite. Sole-hole preservation
uses hole freedom of the retained closure and operation head. Binding freedom
holds independently in the tensor components; the application operation has
no prohibited binding label. Its filling isomorphism combines the proved
serial and parallel addressed-filling laws with the primitive-hole filling
isomorphism. The result is exactly the literal application graph.

\lean{Copying/PropositionOperationalContexts}{10}{25}{Proposition12_2.applicationContext}

\lean{Copying/PropositionOperationalContexts}{26}{34}{Proposition12_2.applicationContext_fill}

The store context is a single hole of type $*\Rightarrow\diamond$ followed
by the focus-free, hole-free store suffix. Its unique address is in the left
serial component. Directed serial paths cannot travel backwards from the
suffix into the hole, and the suffix contains no hole; the existing serial
binding-free theorem therefore applies. Serial addressed filling and the
primitive-hole isomorphism identify each filled context with the required
literal fragment followed by $R$.

\lean{Copying/PropositionOperationalContexts}{35}{41}{Proposition12_2.storeContext}

\lean{Copying/PropositionOperationalContexts}{42}{48}{Proposition12_2.storeContext_fill}

The context-equivalence lemma takes a proved contextual equivalence, an
actual certified graph context, and isomorphisms from its filled endpoints
to the desired graphs. Concrete contextual congruence compares the fillings;
concrete isomorphism invariance compares their endpoint representations.
Reverse the left isomorphism comparison and compose the three equivalences.
The execution laws invoked here are previously proved concrete-machine
results, not premises on this application.

\lean{Copying/PropositionOperationalContexts}{49}{57}{Proposition12_2.context_equiv}

\Needspace{12\baselineskip}
\subsection{Focus-free and hole-free program certificates}

The body is focus free and hole free by the original term-translation
invariants. The closure inherits each property through its lambda edge and
original box. A generic label lemma for the store suffix reduces to four
primitive cases: atom weakening, atom contraction, store and numeral one.
Serial label closure and the canonical-tree edge cases prove it. Specializing
the predicate excludes focus and hole labels at every depth.

\lean{Copying/PropositionAllocationPrograms}{8}{8}{Proposition12_2.body_focusFree}

\lean{Copying/PropositionAllocationPrograms}{9}{10}{Proposition12_2.body_holeFree}

\lean{Copying/PropositionAllocationPrograms}{11}{12}{Proposition12_2.closure_focusFree}

\lean{Copying/PropositionAllocationPrograms}{13}{15}{Proposition12_2.closure_holeFree}

\lean{Copying/PropositionAllocationPrograms}{16}{30}{Proposition12_2.storeSuffix_labels}

\lean{Copying/PropositionAllocationPrograms}{31}{33}{Proposition12_2.storeSuffix_focusFree}

\lean{Copying/PropositionAllocationPrograms}{34}{37}{Proposition12_2.storeSuffix_holeFree}

The two reference programs package the actual local allocation endpoints
with their focus/hole proofs: the source uses the reference operation and
stable numeral, and the target uses the allocation preservation lemmas.
The two beta programs package the original before/after pair with the
unchanged body and stable primitive instance. The existing beta safety theorem
proves all four required endpoint properties. Their interface is
$*\Rightarrow\diamond$, exposing the reference output to the store context.

\lean{Copying/PropositionAllocationPrograms}{38}{44}{Proposition12_2.referenceSource}

\lean{Copying/PropositionAllocationPrograms}{45}{49}{Proposition12_2.referenceTarget}

\lean{Copying/PropositionAllocationPrograms}{50}{54}{Proposition12_2.betaLeft}

\lean{Copying/PropositionAllocationPrograms}{55}{59}{Proposition12_2.betaRight}

The original-program record has the exact Proposition~12.2 left graph and
its translation invariants. The allocated record uses the application label
lemmas with the certified allocated argument. The beta-source and opened
records combine the corresponding certified beta fragment with the already
certified store suffix. These records establish graph admissibility only;
no record field assumes contextual equivalence.

\lean{Copying/PropositionAllocationPrograms}{60}{64}{Proposition12_2.originalProgram}

\lean{Copying/PropositionAllocationPrograms}{65}{71}{Proposition12_2.allocatedProgram}

\lean{Copying/PropositionAllocationPrograms}{72}{76}{Proposition12_2.betaSourceProgram}

\lean{Copying/PropositionAllocationPrograms}{77}{80}{Proposition12_2.openedProgram}

\Needspace{12\baselineskip}
\subsection{Composing the two proved operational links}

The reference pair instantiates the original reference-creation template
with the numeral's proved stability. The beta pair instantiates micro-beta
with the body's proved focus freedom and the primitive instance's proved
stability. Both use identity endpoint isomorphisms. The two directed
contextual improvement theorems apply the already proved robustness endpoints
without forgetting costs.

\lean{Copying/PropositionAllocationEquivalence}{8}{10}{Proposition12_2.referencePair}

\lean{Copying/PropositionAllocationEquivalence}{11}{14}{Proposition12_2.betaPair}

\lean{Copying/PropositionAllocationEquivalence}{15}{17}{Proposition12_2.reference_improvement}

\lean{Copying/PropositionAllocationEquivalence}{18}{21}{Proposition12_2.beta_improvement}

Allocation in the actual original program follows from the reference
pair's termination equivalence in the certified application context. The
left filling is identified with the original program by the source
isomorphism, and the right filling is the allocated graph. The next link is
the allocated-to-beta-source representation isomorphism. Applying the beta
pair in the certified store context proves the third link. Transitivity
then connects the original program to the actual opened graph in every
binding-free closing context, and the observation equivalence characterization
gives the corresponding successful-termination biconditional.

\lean{Copying/PropositionAllocationEquivalence}{22}{29}{Proposition12_2.original_allocated}

\lean{Copying/PropositionAllocationEquivalence}{30}{34}{Proposition12_2.allocated_betaSource}

\lean{Copying/PropositionAllocationEquivalence}{35}{42}{Proposition12_2.betaSource_opened}

\lean{Copying/PropositionAllocationEquivalence}{43}{46}{Proposition12_2.original_opened}

\lean{Copying/PropositionAllocationEquivalence}{47}{50}{Proposition12_2.original_opened_termination}

\Needspace{12\baselineskip}
\subsection{Arbitrary finite and named environments}

\termdef{ambient-weakening-context}
Extend any certified graph context by one genuine primitive weakening using
the existing ordered insertion boundary and its interface equalities. The
old address is lifted through the left parallel component and cast with the
host interface. The new weakening contributes no holes, so the address stays
sole; independent parallel components preserve binding freedom. The filling
isomorphism follows from the parallel addressed-filling law, then the same
interface cast. Its target is exactly outside weakening of the original
filled graph.

\lean{Copying/AmbientWeakeningContext}{10}{25}{Weakening.extendContext}

\lean{Copying/AmbientWeakeningContext}{26}{39}{Weakening.extendContext_fill}

Start with the original primitive closed-term-hole context. Iterating term
extension $k$ times gives the variable context; iterating atom extension $h$
times gives the ambient context. Induction constructs filling isomorphisms
to the previously defined $V_k(G)$ and $A_h(V_k(G))$. The base is the actual
primitive-hole filling theorem; each successor composes the one-extension
filling isomorphism with extension congruence of the induction hypothesis.
Thus the ambient context has its own proved address, sole-hole certificate,
binding freedom and literal filling law.

\lean{Copying/AmbientWeakeningContext}{40}{45}{Weakening.emptyContext}

\lean{Copying/AmbientWeakeningContext}{46}{49}{Weakening.variableContext}

\lean{Copying/AmbientWeakeningContext}{50}{53}{Weakening.ambientContext}

\lean{Copying/AmbientWeakeningContext}{54}{59}{Weakening.variableContext_fill}

\lean{Copying/AmbientWeakeningContext}{60}{66}{Weakening.ambientContext_fill}

A primitive-label predicate survives ambient extension if it holds of the
original graph and both weakening sorts. Induct first on the variable count,
then on the atom count, using the one-extension label equivalence. The
ambient program record specializes this lemma to focus and hole freedom.
Any closed contextual equivalence then lifts to arbitrary finite ambient
environments through the certified ambient context and its filling
isomorphisms. This proves the missing contextual transport of a closed chain,
rather than merely a pair of separate translation-weakening calculations.

\lean{Copying/AmbientWeakeningContext}{67}{77}{Weakening.ambient_allLabels}

\lean{Copying/AmbientWeakeningContext}{78}{84}{Weakening.ambientProgram}

\lean{Copying/AmbientWeakeningContext}{85}{89}{Weakening.ambient_equiv}

The original-in-environment record is precisely the unchanged private-cell
translation at $k,h$; the opened-in-environment record is ambient extension
of $O$. The previous full term-weakening theorem first gives an exact-count
equivalence from the former to ambient extension of the closed original
program. Only then is its cost relation weakened to the universal relation.
Compose with ambient transport of the newly proved closed allocation/beta
chain. The named-environment theorem substitutes the two original list
lengths, and the final observation theorem gives successful-termination
equivalence for every binding-free closing context.

\lean{Copying/PropositionEnvironmentOpening}{10}{14}{Proposition12_2.originalInEnvironment}

\lean{Copying/PropositionEnvironmentOpening}{15}{18}{Proposition12_2.openedInEnvironment}

\lean{Copying/PropositionEnvironmentOpening}{19}{24}{Proposition12_2.environment_weakening_exact}

\lean{Copying/PropositionEnvironmentOpening}{25}{33}{Proposition12_2.original_opened_environment}

\lean{Copying/PropositionEnvironmentOpening}{34}{41}{Proposition12_2.original_opened_named}

\lean{Copying/PropositionEnvironmentOpening}{42}{45}{Proposition12_2.original_opened_environment_termination}

\Needspace{12\baselineskip}
\subsection{Store resolution and concrete regression checks}

Four explicit edge coordinates in $O$ identify the instance retained by
beta, the retained atom contraction, the retained store and numeral one.
Each label equation is definitional: nested serial tags select the original
edge without relabelling it. The instance target and a contraction source
share the outer atom junction; the contraction target and store source
share the inner atom junction. These two concrete incidence witnesses prove
the two atom hops.

\lean{Copying/PropositionOpeningChecks}{8}{9}{Proposition12_2.openedInstance}

\lean{Copying/PropositionOpeningChecks}{10}{10}{Proposition12_2.openedContraction}

\lean{Copying/PropositionOpeningChecks}{11}{11}{Proposition12_2.openedStore}

\lean{Copying/PropositionOpeningChecks}{12}{13}{Proposition12_2.openedOne}

\lean{Copying/PropositionOpeningChecks}{14}{14}{Proposition12_2.opened_instance_label}

\lean{Copying/PropositionOpeningChecks}{15}{15}{Proposition12_2.opened_contraction_label}

\lean{Copying/PropositionOpeningChecks}{16}{16}{Proposition12_2.opened_store_label}

\lean{Copying/PropositionOpeningChecks}{17}{19}{Proposition12_2.opened_one_label}

\lean{Copying/PropositionOpeningChecks}{20}{26}{Proposition12_2.opened_instance_hop}

\lean{Copying/PropositionOpeningChecks}{27}{31}{Proposition12_2.opened_contraction_hop}

Composing the two hops and the store-label equation proves resolution in
exactly two hops. The already proved store-trace uniqueness theorem shows
that every successful resolution from that instance has this same store and
length. A separate literal junction witnesses adjacency of the store to the
numeral-one edge. The allocation non-stability theorem confirms that the
whole heap argument cannot satisfy beta's stable-fragment premise.

Finite enumeration checks the outer edge counts: five in the original graph,
eight after allocation, eight in the reassociated beta source, and nine after
opening the box. The increase after beta is compatible with opening previously
nested body edges; no claim of a graph-size decreasing rule is made. Finally,
the environment theorem is instantiated at two variables and three atoms.
Termination equivalence is quantified over all closing observations, so no
assumption of terminating callbacks or single-use closures is introduced.

\lean{Copying/PropositionOpeningChecks}{32}{34}{Proposition12_2.opened_resolves}

\lean{Copying/PropositionOpeningChecks}{35}{38}{Proposition12_2.opened_resolves_unique}

\lean{Copying/PropositionOpeningChecks}{39}{43}{Proposition12_2.opened_store_content}

\lean{Copying/PropositionOpeningChecks}{44}{46}{Proposition12_2.allocated_argument_not_stable}

\lean{Copying/PropositionOpeningChecks}{47}{52}{Proposition12_2.intermediate_edge_counts}

\lean{Copying/PropositionOpeningChecks}{53}{55}{Proposition12_2.mixed_environment_opening}

\paragraph{Dependency map and remaining work.}
Completed: empty-forest representation isomorphisms and original padding
$\to$ exact source normalization; branch/serial reassociation $\to$ stable
instance/store separation; proved reference and beta template consequences
plus actual binding-free contexts $\to$ closed original-to-opened equivalence;
full weakening plus a certified ambient context $\to$ arbitrary finite and
named-environment equivalence. Concrete store resolution is proved on the
opened endpoint. No semantic leaf is an assumed template or machine theorem.

Still unfinished: normalize the retained translated body and sharing/box
structure into the required parametricity instance; apply parametricity and
normalize to the unchanged constant-result translation; complete the chain
at the required common depth allowances and discharge the remaining
correspondence/substitution obligations. Proposition~12.2 itself is not yet
proved. The estimate remains 70--80\% of mathematical effort, with low
confidence: this is a concrete program-equivalence milestone, while the
remaining structural normalization and correspondence work are still hard
to size.
